\section{Discussion}
\label{sec:discussion}

\paragraph{What the farm test measures.} Every model we trained, including the two that earlier work recommended,
reaches 98.5 to 99.6\% on the random farm test and loses 21 to 25 points on photographs from a second collection of
the same group. The original authors did not expect distribution shift to matter~\cite{machuve2022}. Our results show
that it matters more than the choice of network. We cannot say from these data which difference between the
collections causes the drop, but the farm collection itself suggests a strong candidate. Its salmonella and
Newcastle photographs come from experimentally inoculated birds and its healthy and coccidiosis photographs from
farms~\cite{machuve2021farm}, while every PCR-labelled photograph comes from a farm. A model can then learn the
setting of each class, such as the surface, the lighting or the birds, along with the dropping, which is the kind of
shortcut described by Geirhos et al.~\cite{geirhos2020}. Healthy droppings in the farm collection also fill a median
of 2.6\% of the photograph, against 12 to 22\% for the other classes, so framing could act as a shortcut too. Three
observations fit this reading without proving it. Exact duplicates are not the cause, because removing them leaves
the farm score unchanged, and neither are the conflicting PCR files, which move every model by about one point. Most
errors are missed detections and calls of healthy or salmonella-like, not swaps between the diseases. And on
normal-looking PCR-labelled photographs the models lean towards salmonella (Section~\ref{sec:grades-results}), the class
whose training photographs came from one inoculation setting. We have not tested these shortcuts directly. Sakib et
al.\ saw a much smaller drop, from 99.19 to 95.62\%, on a local-farm set~\cite{sakib2026}. Their abstract does not say
how that set was labelled or collected, so we cannot tell whether it shares the setting of the training photographs.
Bodempudi et al.\ reached at most 85\% across regions for a binary \emph{Salmonella} task~\cite{bodempudi2026}.
Together these results suggest that published accuracies on this task should be read as within-collection
accuracies.

\paragraph{Why use a detector.} On photographs our detector is not more accurate than a whole-image classifier. Its
value is that it says where the dropping is. That is what makes a fixed pen camera usable, because under our size
assumption a dropping covers at most about 0.1\% of a pen frame, and it is also what lets the explanation step
measure the colour of the dropping itself rather than the litter around it. The pen-scale fine-tune shows that the
detector can learn the size and setting of a pen camera from public video and pasted droppings alone, without a
single labelled pen dropping. It lost 0.02 mAP50 on close-ups in the process. The deployed Pi pipeline does not use
the pen model yet. It finds candidates with a still-spot finder and judges them with the close-up model, whose false
alarms on white birds the explanation step filters out.

\paragraph{Why a microkernel.} A monitor in a barn runs unattended, and a crash in a camera driver or a model runtime
should cost a few frames, not the device. On QNX each stage of our pipeline is an ordinary process, and the
supervisor relaunched any stage we killed within 20 to 33\,ms of noticing the exit while the others kept running. The
same structure is possible on Linux with process supervision. What QNX adds is that the drivers, the file system and
the network stack are also user processes that can fail and restart the same way~\cite{qnx8sysarch}, but our tests
killed only our own stages, not a QNX driver or service. We did not compare fault recovery against Linux on the same
board, so we do not claim QNX recovers faster. Running the model on QNX also needed one fix in the platform's
packages. The Python binding of ONNX Runtime in the QNX open-source repository does not load, because its build
script treats QNX as Windows and looks for the wrong extension module. We worked around it with a 111-line C bridge
to the runtime's C library and have prepared a patch for the package recipe.

\paragraph{Limitations.} All photographs come from one region of Tanzania and one research group. No farm photograph
carries a laboratory diagnosis of its own dropping, two farm classes come from inoculated birds, and the PCR release
contains files filed under two classes, so we treat the classes as screening categories, not diagnoses. The farm
split is random by photograph because the release carries no farm or visit information. The pen test and the test
clip use real droppings pasted onto real barn frames, so they measure whether a pen-sized dropping can be found on
unseen litter, not field accuracy on naturally occurring droppings. The pen size of 8 to 30 pixels is an assumption.
We have no labelled droppings from a real pen camera, and we have not measured the detection accuracy of the deployed
pipeline. The detector and each baseline were trained once, with one seed and one set of hyperparameters that we did
not tune, and EfficientNet-B0 for only 8 epochs. Run-to-run variation is therefore not measured. To explain the drop it
would have to exceed 20 points, far beyond the 0.02 mAP50 gap between our two detector runs on the farm test.
Finally, the edge timings are for inference only and do not include image decoding, post-processing or power draw.
The appearance grades come from three AI models following a rubric, not from veterinarians. Their agreement shows
the grades are consistent, not that they are right, and three models from one vendor may share blind spots. The
rubric is new and has not been validated against expert grading.

\paragraph{Future work.} The most useful next step is labelled droppings from fixed pen cameras on more than one
farm, split by farm. Short of that, the farm collection could be re-split with near-duplicate removal before
splitting, the size shortcut could be tested by cropping every class to the same framing, and the setting shortcut by
testing on salmonella and Newcastle photographs taken on farms. On the device, the next steps are a Linux baseline on
the same Raspberry Pi 5 and fault tests that kill QNX services, not only our own stages. The box score needs
calibration before it is shown to a farmer, and the 0.25 threshold should be chosen from a risk-coverage curve on
validation data rather than fixed.

\section{Conclusion}
\label{sec:conclusion}

Classifiers of poultry droppings that score 98.5 to 99.6\% on a random split of the Tanzanian farm photographs score
75 to 78\% on PCR-labelled photographs from a second collection taken on farms, and this holds for every architecture
we tried, including our detector. We therefore recommend that future work on this task report a second collection
alongside the farm test. A detector cannot fix this drop, but it is what makes a fixed pen camera usable. Fine-tuned
on public barn video with pasted droppings, our YOLO11n detector finds 80\% of pen-sized droppings on unseen cameras
with 93\% precision. A pipeline built around the close-up detector runs on a Raspberry Pi 5 under QNX, where a stage
killed with \texttt{SIGKILL} is relaunched within tens of milliseconds while the rest keeps running.
